\section{Experiments}
\label{sec:experiments}

前文已经分别从问题定义、数据协议与方法实现层面对 IDC patch classification 的实验前提进行了形式化刻画，因此本节不再重复训练流水线或配置项本身，而是集中报告经验结果。具体而言，本节围绕两个实验问题展开：其一，在统一 patient-wise 协议下，不同 ResNet 主干网络（backbone）的性能表现有何差异；其二，当数据划分策略由 patient-wise split 改为 patch-wise random split 时，训练集与验证集上的指标会呈现出怎样的变化。

\subsection{Backbone Baseline Comparison}

我们首先考察在严格 patient-wise 设定下，不同 ResNet 主干网络所形成的基线差异。为保证比较具有可解释性，本文仅替换主干结构本身，而保持其余训练条件不变，包括输入尺寸、数据预处理、优化超参数、训练轮数、正类权重（positive-class weight）与决策阈值。

\begin{table}[t]
    \centering
    \small
    \setlength{\tabcolsep}{4pt}
    \caption{Patient-wise split 下不同 ResNet backbone 的基线比较。所有模型均在相同训练协议下训练：输入尺寸为 $50 \times 50$，使用 ImageNet 预训练权重，训练 $20$ 个 epoch，batch size 为 $256$，学习率为 $3\times10^{-4}$，正类权重（positive-class weight）为 $2.5$，并以固定阈值 $0.5$ 在同一验证集上评估。最佳结果以粗体标出。}
    \label{tab:backbone_baseline_comparison}
    \begin{tabular*}{\textwidth}{@{\extracolsep{\fill}}lccccc@{}}
        \toprule
        Backbone & Precision & Recall & F1-score & BAcc. & AUC \\
        \midrule
        ResNet-18 & 0.9089 & 0.9616 & 0.9345 & 0.9603 & 0.9922 \\
        ResNet-34 & 0.9069 & 0.9593 & 0.9324 & 0.9587 & 0.9921 \\
        ResNet-50 & \textbf{0.9334} & \textbf{0.9830} & \textbf{0.9575} & \textbf{0.9766} & \textbf{0.9976} \\
        \bottomrule
    \end{tabular*}
\end{table}


从验证结果看，ResNet-18 与 ResNet-34 的整体表现接近，其中 ResNet-34 在 Precision、Recall、F1-score、Balanced Accuracy 与 ROC-AUC 上均略低于 ResNet-18。相比之下，ResNet-50 在各项关键指标上均明显优于前两者，尤其在 F1-score、Balanced Accuracy 与 ROC-AUC 上给出当前最优结果。

基于当前结果，patient-wise 设定下的 backbone 表现并未随网络深度增加而呈现简单单调趋势；在本次实验中，ResNet-50 是三者中验证指标最优的 backbone。与此同时，\cref{tab:backbone_baseline_comparison} 报告的是统一协议、固定随机种子与单次正式运行下的结果，可作为后续比较的主基线。

\subsection{Protocol Comparison: Patient-wise vs. Patch-wise}

在完成 backbone 基线比较之后，我们进一步比较 patient-wise split 与 patch-wise random split 两种数据划分方式。为保证可比性，实验在统一训练配置下进行，并同时报告训练集与验证集指标以及相应的训练-验证差值（train-eval gap）。

\begin{table}[t]
    \centering
    \footnotesize
    \setlength{\tabcolsep}{2pt}
    \caption{不同数据划分协议（split protocol）下的训练集与验证集落差比较。除切分策略外，其余训练配置保持一致。这里重点报告 F1-score、Balanced Accuracy 与 ROC-AUC 在训练集与验证集之间的泛化落差（generalization gap），以观察 patch-wise random split 是否更容易产生“训练很高、验证明显回落”的虚高现象。}
    \label{tab:split_protocol_comparison}
    \begin{tabular*}{\textwidth}{@{\extracolsep{\fill}}p{1.45cm}lcccccc@{}}
        \toprule
        Split & Model & Train F1 & Eval F1 & F1 Gap & BAcc. Gap & AUC Gap & Eval AUC \\
        \midrule
        Patch & ResNet-18 & 0.9553 & 0.9103 & 0.0450 & 0.0294 & 0.0103 & 0.9863 \\
        Patient & ResNet-18 & 0.9573 & 0.9345 & 0.0228 & 0.0156 & 0.0048 & 0.9922 \\
        Patch & ResNet-34 & 0.9547 & 0.9138 & 0.0408 & 0.0246 & 0.0093 & 0.9874 \\
        Patient & ResNet-34 & 0.9594 & 0.9324 & 0.0271 & 0.0185 & 0.0052 & 0.9921 \\
        Patch & ResNet-50 & 0.9731 & 0.9381 & 0.0350 & 0.0205 & 0.0072 & 0.9915 \\
        Patient & ResNet-50 & 0.9744 & 0.9575 & 0.0169 & 0.0089 & 0.0013 & 0.9976 \\
        \bottomrule
    \end{tabular*}
\end{table}


从结果上看，在 ResNet-18、ResNet-34 与 ResNet-50 三种 backbone 上，patch-wise split 的训练指标均保持在较高水平，但其训练-验证差值始终大于 patient-wise split。以 F1-score 为例，patch-wise 的 train-eval gap 分别为 $0.0450$、$0.0408$ 与 $0.0350$，而 patient-wise 对应为 $0.0228$、$0.0271$ 与 $0.0169$；在 Balanced Accuracy 上，patch-wise 的 gap 分别为 $0.0294$、$0.0246$ 与 $0.0205$，patient-wise 则为 $0.0156$、$0.0185$ 与 $0.0089$；在 ROC-AUC 上，patch-wise 的 gap 分别为 $0.0103$、$0.0093$ 与 $0.0072$，而 patient-wise 分别为 $0.0048$、$0.0052$ 与 $0.0013$。
\begin{itemize}
    \item 在三种 backbone 上，patch-wise split 的训练指标整体更高，但对应的训练-验证差值也始终更大。
    \item 在 F1-score 上，patch-wise 的 train-eval gap 分别为 $0.0450$、$0.0408$ 与 $0.0350$，明显高于 patient-wise 的 $0.0228$、$0.0271$ 与 $0.0169$。
    \item 在 Balanced Accuracy 上，patch-wise 的 gap 分别为 $0.0294$、$0.0246$ 与 $0.0205$，高于 patient-wise 的 $0.0156$、$0.0185$ 与 $0.0089$。
    \item 在 ROC-AUC 上，patch-wise 的 gap 分别为 $0.0103$、$0.0093$ 与 $0.0072$，而 patient-wise 分别为 $0.0048$、$0.0052$ 与 $0.0013$。
    \item 整体上，patch-wise random split 更容易获得较高的训练分数，同时也伴随着更大的训练-验证差值；相比之下，patient-wise split 的训练与验证指标关系更为接近。
\end{itemize}

表中的现象为后续讨论实验协议带来的评估差异提供了直接经验基础。
